\chapter*{Abstract}
\addcontentsline{toc}{chapter}{Abstract}

High-density Single Nucleotide Polymorphism (SNP) arrays provide
genome-wide information on genetic variation and have been
successfully used for breed discrimination and assignment in
livestock populations
\cite{dimauro2013canonical, dimauro2015sheep,
      sorbolini2016signatures}.
However, their high dimensionality introduces substantial
computational challenges and raises the question of whether the full
marker set is necessary for reliable breed assignment.

This work investigates machine learning and deep representation
learning for livestock breed classification from high-dimensional
genotype data. Particular attention is given to two practical
questions: how breed-classification performance on unseen individuals
changes as the number of training individuals per breed is reduced,
and how many genomic markers are needed to retain most of the
predictive information available in the original SNP array.

Generative and contrastive representation-learning approaches are
combined with feature-attribution techniques to derive global
rankings of informative SNPs. The resulting reduced marker panels are
compared with established population-genetics and machine-learning
ranking strategies, including fixation-index-based ranking and Random
Forest feature importance
\cite{weir1984estimating, breiman2001random}.

The experimental protocol evaluates breed-classification performance
under progressively reduced training-set sizes and SNP-panel sizes.
The primary analysis is conducted on ovine genomic data, and the same
methodological protocol is subsequently applied to an independent
caprine dataset to assess whether the observed computational behavior
generalizes across livestock species.

Overall, the study investigates whether compact genomic
representations and reduced SNP panels can preserve reliable breed
classification while decreasing data and marker requirements. The
resulting framework is designed to quantify the trade-off between
predictive performance, training sample size, and genomic panel
dimensionality.
